\chapter{TypeScript e la gestione delle dipendenze}

\section{Carta d'identit\`a di TypeScript}

\begin{esamebox}
\emph{``Descrivere le caratteristiche del linguaggio TypeScript in quanto estensione di
JavaScript (le caratteristiche proprie, non quelle ereditate da JavaScript).''}
\end{esamebox}

\begin{definizione}[TypeScript]
\textbf{TypeScript} aumenta la sintassi di JavaScript con un \textbf{sistema di tipi}:
\begin{itemize}
  \item si pu\`o dichiarare il \emph{tipo atteso} di variabili, parametri e valori di ritorno
        delle funzioni;
  \item \`e un linguaggio \textbf{compilato} (traspilato): i programmi TS vengono tradotti in
        programmi JS. Di conseguenza i costrutti aggiuntivi di TS sono \textbf{esclusivamente
        dichiarativi e non operazionali}: ``scompaiono'' con la compilazione e non hanno alcun
        impatto sull'esecuzione;
  \item implementa un sistema di \textbf{inferenza dei tipi}: a compile-time determina il tipo
        risultante di funzioni ed espressioni e verifica la consistenza con i tipi
        ``promessi'' o con l'uso dei risultati;
  \item d\`a il meglio di s\'e \textbf{integrato in un IDE}: l'inferenza viene applicata
        gi\`a mentre si scrive, prevenendo molti errori prima ancora della compilazione.
\end{itemize}
\end{definizione}

\begin{importante}[Tipi statici vs dinamici]
In linguaggi come JavaScript o C il tipo influenza lo \emph{slot di memoria} usato per il
dato; in TypeScript no: le dichiarazioni di tipo \textbf{non hanno effetto a run-time},
perch\'e svaniscono nella compilazione. A cosa servono allora? A permettere, \emph{a
compile-time}, la verifica della consistenza dei tipi in espressioni e assegnamenti: il
programmatore esprime le proprie \emph{intenzioni} e il compilatore (e l'IDE) rileva le
discrepanze fra intenzione e codice scritto. Permettono inoltre di controllare che
propriet\`a e metodi a cui si accede \emph{esistano davvero}, e di introdurre nelle classi i
\emph{modificatori di accesso} (public, private\dots) assenti in JavaScript.
\end{importante}

\begin{attenzione}[Il tipo any]
Si pu\`o ``non esprimere'' il tipo di un'entit\`a dichiarandolo \texttt{any} (qualsiasi cosa),
ma cos\`i si perdono buona parte dei vantaggi di TypeScript: da evitare, salvo i rari casi in
cui si deve far convivere codice TS con codice JavaScript privo di dichiarazioni.
\end{attenzione}

\section{Package management: npm}

Con TypeScript il progetto inizia a dipendere dal compilatore \texttt{tsc} e dai
\emph{package} che definiscono il linguaggio. \`E un'istanza di un caso generale: nello
sviluppo web \`e comune usare package di terze parti, distribuiti tramite \emph{repository
online}. Domande tipiche: quando ridistribuire il progetto perch\'e lo richiedono le
dipendenze? Come evitare che l'aggiornamento di un package ``rompa'' applicazioni gi\`a
distribuite? Che succede se uno sviluppatore ritira il proprio package?

\begin{definizione}[Package manager]
Un \textbf{package manager} \`e un tool (tipicamente da linea di comando) che permette di
accedere al registry/repository abbinato, scaricare e installare i package necessari e
\emph{dichiararne la dipendenza} tramite opportuni dependency file. Il pi\`u diffuso nel
mondo JS/TS \`e \textbf{npm} (node package manager), col registry \texttt{npmjs.com};
il modo pi\`u semplice di installarlo \`e installare \textbf{Node.js}, che lo include.
\end{definizione}

\begin{lstlisting}
> npm install [-D] [-g] nome-package[@versione]
\end{lstlisting}
\begin{itemize}
  \item \texttt{-D}: \emph{dev dependency} (default: runtime dependency);
  \item \texttt{-g}: installazione \emph{globale} (default: locale al progetto);
  \item \texttt{@versione}: quale versione installare; \texttt{@latest} = ultima stabile,
        \texttt{@next} = prossima uscita (non stabile).
\end{itemize}

\begin{importante}[Dipendenze: locali vs globali, dev vs runtime]
\textbf{Locali vs globali}: npm pu\`o installare i package a livello globale (come librerie
di sistema) o \emph{locale}, nella cartella del progetto. La strategia pi\`u comune oggi \`e
l'installazione \textbf{locale}: rende lo sviluppo indipendente dalla configurazione del
dispositivo e facilita la collaborazione via Git/GitLab.
\textbf{Dev vs runtime}:
\begin{itemize}
  \item \emph{dev dependencies}: servono in fase di sviluppo ma ``scompaiono'' alla
        distribuzione -- TypeScript \`e una di queste (compilato in JS, compilatore e
        librerie TS non servono pi\`u);
  \item \emph{runtime dependencies}: codice effettivamente invocato durante l'esecuzione, da
        includere nella distribuzione.
\end{itemize}
La distinzione \`e rilevante nel ``confezionamento'' dei progetti: sviluppo leggibile,
distribuzione concisa ed efficiente.
\end{importante}

\subsection{Struttura di progetto npm}
\begin{itemize}
  \item \textbf{node\_modules}: il codice dei package; gestita interamente da npm, non va
        modificata e \textbf{va esclusa da Git};
  \item \textbf{package.json}: configurazione del progetto e dipendenze; npm ci scrive, ma \`e
        modificabile anche a mano
        (es. \texttt{"devDependencies": \{"typescript": "\^{}5.6.3"\}});
  \item \textbf{package-lock.json}: dipendenze e sotto-dipendenze; generato da npm, da non
        modificare.
\end{itemize}

\section{Configurare e compilare TypeScript}

Il compilatore \textbf{tsc} \`e uno script JavaScript contenuto nel package typescript,
eseguito tramite il tool \textbf{npx} (incluso in npm).
\begin{lstlisting}
> npx tsc --init    # crea tsconfig.json con impostazioni predefinite
> npx tsc           # compila: cerca i .ts e salva i .js nella outDir
\end{lstlisting}

Le \textbf{compilerOptions} principali in \texttt{tsconfig.json}:
\begin{itemize}
  \item \texttt{target}: versione di JavaScript dell'output. Default \texttt{es2016}
        (eseguibile da tutti i browser); nel corso si usa \texttt{es2020} per avere pi\`u
        feature (ormai ampiamente supportato);
  \item \texttt{module}: come gestire i moduli. Prima che ECMAScript introducesse
        import/export, Node.js aveva una propria sintassi (``CommonJS''); poich\'e scriviamo
        per il browser e vogliamo la sintassi ES, si indica \texttt{es2015} (la versione che
        ha introdotto import/export); per Node moderno si userebbe \texttt{nodenext};
  \item \texttt{esModuleInterop}: interoperabilit\`a fra le due modalit\`a di specifica dei
        moduli (rilevante soprattutto programmando per Node con stile ES);
  \item \texttt{forceConsistentCasingInFileNames}: il compilatore segnala import di uno
        stesso modulo con uso inconsistente di maiuscole/minuscole (problema fra S.O. che
        distinguono e S.O. che no);
  \item \texttt{strict}: maggiore severit\`a nella verifica dei tipi;
  \item \texttt{skipLibCheck}: non verifica i tipi dentro le librerie;
  \item \texttt{outDir}: \emph{dove} mettere i file compilati. Conviene tenere i sorgenti
        \texttt{.ts} in una cartella \texttt{src} e compilare in \texttt{scripts}: le pagine
        HTML si riferiscono a \texttt{scripts} (la ``vera'' applicazione da distribuire) e
        ignorano i sorgenti; viceversa ai collaboratori serve solo \texttt{src}, perch\'e
        \texttt{scripts} si rigenera compilando. tsc replica in output la struttura di
        cartelle dei sorgenti.
\end{itemize}

\section{Dichiarazioni di tipo}

\begin{lstlisting}[language=JavaScript]
let nomevar: TIPO;
const c: TIPO = espressione;
function fun(p: TIPO, q: TIPO): TIPORET {...}
const arr = (a: TIPO, b: TIPO): TIPORET => ...
\end{lstlisting}
Le dichiarazioni di tipo di \emph{variabili e parametri} sono obbligatorie, a meno che non
venga subito assegnato un valore (TS inferisce il tipo dall'espressione). Il tipo di
\emph{ritorno} non \`e obbligatorio (inferito dalle \texttt{return}), ma dichiararlo
esplicitamente \`e utile: esprime le intenzioni, e se il valore restituito \`e del tipo
sbagliato -- o manca la return -- il compilatore lo segnala.

\section{Espressioni di tipo}

Imparare TypeScript significa soprattutto saper scrivere le \textbf{espressioni di tipo} --
meno banale di quanto sembri (``una funzione che prende un array di interi e restituisce una
stringa''\dots). Scriverle correttamente \`e una forma preliminare di \emph{testing}: se il
codice supera i controlli di tipo, molti errori grossolani sono gi\`a esclusi.

\subsection{Tipi base, array, tuple}
\begin{itemize}
  \item nomi dei tipi base: \texttt{string}, \texttt{boolean}, \texttt{number},
        \texttt{undefined}, \texttt{null}, \texttt{bigint}, \texttt{symbol};
  \item il nome di una \textbf{classe}: la variabile conterr\`a solo oggetti di quella classe;
  \item \texttt{void}: tipo di ritorno di una funzione che non ritorna nulla;
  \item \textbf{array}: tipo degli elementi + quadre --
        \texttt{let mioArray: number[];} \texttt{let iscritti: Studente[];}
  \item \textbf{tupla}: array di lunghezza fissa con tipi posizionali noti --
        \texttt{let tripla: [number, string, Studente]}.
\end{itemize}

\subsection{Tipo unione}
Composizione con \texttt{|}: un valore ``di tipo A \emph{oppure} B''.
\begin{lstlisting}[language=JavaScript]
let misto: (number | string)[]
function trovaStudente(matricola: number): Studente | null {...}
// unione di tipi literal (elenco di valori possibili):
let cifra: 0|1|2|3|4|5|6|7|8|9;
let primavera: "marzo" | "aprile" | "maggio";
\end{lstlisting}

\subsection{Tipo funzione}
Si costruisce da tipi dei parametri e del ritorno:
\texttt{(p1: T1, \dots, pN: TN) => TRET}. I ``nomi'' \texttt{p1\dots pN} sono solo
segnaposto.
\begin{esempio}
\texttt{superCalcolo} lancia un calcolo remoto e prende una callback che ricever\`a il
risultato (array di interi) restituendo un boolean:
\begin{lstlisting}[language=JavaScript]
function superCalcolo(datiInput: number[],
    callback: (datiOutput: number[]) => boolean): void {...}
// e per memorizzare superCalcolo in una variabile:
let somevar: (x: number[], y: (z: number[]) => boolean) => void;
somevar = superCalcolo;
\end{lstlisting}
\end{esempio}

\subsection{Le classi in TypeScript}
Una classe TS \emph{dichiara le propriet\`a} (variabili di istanza) degli oggetti, con i tipi:
\begin{lstlisting}[language=JavaScript]
class Studente {
  nome: string;
  matricola: number;
  votiEsami: number[];
  ...
}
\end{lstlisting}
Costruttore e metodi come in JS, ma con i tipi dichiarati. Inoltre:
\begin{itemize}
  \item \textbf{visibilit\`a}: \texttt{private}, \texttt{public} (default),
        \texttt{protected}; in pi\`u \texttt{readonly} (assegnata nel costruttore e mai
        pi\`u modificata);
  \item TS verifica che tutte le variabili d'istanza siano \emph{inizializzate nel
        costruttore}; per saltare il controllo (perch\'e verr\`a inizializzata dopo, in
        conseguenza di un evento) si pospone \texttt{!}; per dichiararla \emph{opzionale}
        (potrebbe non avere valore) si pospone \texttt{?};
  \item propriet\`a e metodi devono \emph{sempre rispettare la dichiarazione della classe}: le
        uniche propriet\`a aggiungibili/rimovibili durante la vita dell'oggetto sono quelle
        opzionali (\texttt{?}).
\end{itemize}

\subsection{Le interface}
\begin{definizione}[Interface]
Le \textbf{interface} sono un costrutto \emph{completamente dichiarativo}, senza semantica
operazionale: non se ne trova traccia nel codice compilato. Servono a garantire che gli
oggetti in input/output abbiano determinate propriet\`a -- importante in un linguaggio dove si
creano spesso oggetti \emph{senza} costruttori. Un'interface \`e un elenco di propriet\`a
(anche propriet\`a-funzione); il suo nome definisce un \emph{tipo}; con \texttt{?} si
dichiarano propriet\`a opzionali.
\begin{lstlisting}[language=JavaScript]
interface Book {
  title: string;
  author: string;
  year: number
}
\end{lstlisting}
\end{definizione}

Quando un oggetto \emph{rispetta} un'interfaccia \texttt{I}?
\begin{itemize}
  \item un \textbf{object literal} la rispetta se \emph{corrisponde esattamente}: tutte le
        propriet\`a non opzionali, e \emph{nessuna} propriet\`a aggiuntiva;
  \item un oggetto di una \textbf{classe C} la rispetta se C dichiara tutte le propriet\`a
        non opzionali di I; C pu\`o averne anche altre. Una classe che definisce tutte le
        propriet\`a di I si dice che \textbf{implementa} l'interfaccia, esplicitabile con
        \texttt{class C implements I \{\dots\}}.
\end{itemize}
La versione TypeScript della DOM API definisce come interface tutti i tipi di oggetti che
usa. Molte librerie nate in JavaScript forniscono definizioni di tipo parallele in file
\textbf{.d.ts} (es. i tipi della DOM API sono in \texttt{lib.dom.d.ts}, nel package
typescript).

\subsection{Dizionari indicizzati}
In JS/TS \`e comune usare gli oggetti come ``dizionari'' (simili a hash map) per reperire
valori da una chiave. Ma TypeScript \emph{non permette} di accedere a propriet\`a non
menzionate dalla definizione di tipo -- salvo che l'oggetto sia dichiarato esplicitamente
come dizionario indicizzato:
\begin{lstlisting}[language=JavaScript]
const repo: {[chiave: string]: Persona} = {};
function memorizza(p: Persona) { repo[p.codiceFiscale] = p; }
\end{lstlisting}
Le chiavi possono essere soltanto \texttt{string} o \texttt{number}; i valori di qualsiasi
tipo.

\subsection{Accessors: getter e setter}
Caratteristica di JavaScript (non di TypeScript), ma la sua utilit\`a emerge con i tipi.
Permettono di definire \textbf{propriet\`a ``calcolate''}: utile ad esempio quando
un'interfaccia richiede \emph{propriet\`a} che per una classe sarebbero naturalmente
\emph{metodi} (calcolabili da altre propriet\`a).
\begin{lstlisting}[language=JavaScript]
interface RagioneSociale { denominazione: string; codiceFiscale: string; }
class Persona {
  get denominazione(): string { return `${this.nome} ${this.cognome}`; }
  set denominazione(den: string) {
    const sep = den.split(" ");
    this.nome = sep[0]; this.cognome = sep[1];
  }
}
tizio.denominazione = "Mickey Mouse"; // invoca il setter
const s = tizio.denominazione;        // invoca il getter
\end{lstlisting}
Regole:
\begin{itemize}
  \item il \texttt{get} non prende parametri e restituisce il valore; il \texttt{set} non
        restituisce nulla e prende il nuovo valore;
  \item se presenti entrambi, il tipo di ritorno del get deve corrispondere al parametro del
        set;
  \item non \`e obbligatorio averli entrambi: solo get $\Rightarrow$ propriet\`a
        \emph{readonly}; solo set $\Rightarrow$ scrivibile ma non leggibile.
\end{itemize}
Usi tipici: propriet\`a calcolate da altri valori; controllo degli accessi senza rendere
private le propriet\`a interne (limitare valori numerici, validare formati, elaborare prima
di restituire).

\subsection{Generics}
\begin{definizione}[Generics]
I \textbf{generics} esprimono un comportamento \emph{indipendente dal tipo di dato} su cui
operano: un ordinamento funziona su array di tipi diversi (restituendo un array dello stesso
tipo); uno Stack ha push/pop analoghi qualunque sia il tipo degli elementi. TypeScript
permette funzioni, classi e interfacce \emph{generiche}, senza doverle ridefinire per ogni
tipo. La sintassi usa le parentesi angolari; con una \emph{variabile di tipo} si creano
elementi generici:
\begin{lstlisting}[language=JavaScript]
const insiemeNum = new Set<number>()
const insiemeStudenti = new Set<Studente>()
function middle<T>(lista: T[]): T {
  return lista[lista.length/2];
}
\end{lstlisting}
(Se una classe \`e generica, lo \`e anche il suo costruttore.)
\end{definizione}
